\documentclass{amsart}
% For dealing with references we use the comment environment
\usepackage{verbatim}
\newenvironment{reference}{\comment}{\endcomment}
%\newenvironment{reference}{}{}
\newenvironment{slogan}{\comment}{\endcomment}
\newenvironment{history}{\comment}{\endcomment}

% For commutative diagrams we use Xy-pic
\usepackage[all]{xy}

% We use 2cell for 2-commutative diagrams.
\xyoption{2cell}
\UseAllTwocells

% We use multicol for the list of chapters between chapters
\usepackage{multicol}

% This is generally recommended for better output
\usepackage{lmodern}
\usepackage[T1]{fontenc}

% For cross-file-references

% Package for hypertext links:
\usepackage{hyperref}

% For any local file, say "hello.tex" you want to link to please

% Theorem environments.
%
\theoremstyle{plain}
\newtheorem{theorem}[subsection]{Theorem}
\newtheorem{proposition}[subsection]{Proposition}
\newtheorem{lemma}[subsection]{Lemma}

\theoremstyle{definition}
\newtheorem{definition}[subsection]{Definition}
\newtheorem{example}[subsection]{Example}
\newtheorem{exercise}[subsection]{Exercise}
\newtheorem{situation}[subsection]{Situation}

\theoremstyle{remark}
\newtheorem{remark}[subsection]{Remark}
\newtheorem{remarks}[subsection]{Remarks}

\numberwithin{equation}{subsection}

% Macros
%
\def\lim{\mathop{\mathrm{lim}}\nolimits}
\def\colim{\mathop{\mathrm{colim}}\nolimits}
\def\Spec{\mathop{\mathrm{Spec}}}
\def\Hom{\mathop{\mathrm{Hom}}\nolimits}
\def\Ext{\mathop{\mathrm{Ext}}\nolimits}
\def\SheafHom{\mathop{\mathcal{H}\!\mathit{om}}\nolimits}
\def\SheafExt{\mathop{\mathcal{E}\!\mathit{xt}}\nolimits}
\def\Sch{\mathit{Sch}}
\def\Mor{\mathop{\mathrm{Mor}}\nolimits}
\def\Ob{\mathop{\mathrm{Ob}}\nolimits}
\def\Sh{\mathop{\mathit{Sh}}\nolimits}
\def\NL{\mathop{N\!L}\nolimits}
\def\CH{\mathop{\mathrm{CH}}\nolimits}
\def\proetale{{pro\text{-}\acute{e}tale}}
\def\etale{{\acute{e}tale}}
\def\QCoh{\mathit{QCoh}}
\def\Ker{\mathop{\mathrm{Ker}}}
\def\Im{\mathop{\mathrm{Im}}}
\def\Coker{\mathop{\mathrm{Coker}}}
\def\Coim{\mathop{\mathrm{Coim}}}

% Boxtimes
%
\DeclareMathSymbol{\boxtimes}{\mathbin}{AMSa}{"02}

%
% Macros for moduli stacks/spaces
%
\def\QCohstack{\mathcal{QC}\!\mathit{oh}}
\def\Cohstack{\mathcal{C}\!\mathit{oh}}
\def\Spacesstack{\mathcal{S}\!\mathit{paces}}
\def\Quotfunctor{\mathrm{Quot}}
\def\Hilbfunctor{\mathrm{Hilb}}
\def\Curvesstack{\mathcal{C}\!\mathit{urves}}
\def\Polarizedstack{\mathcal{P}\!\mathit{olarized}}
\def\Complexesstack{\mathcal{C}\!\mathit{omplexes}}
% \Pic is the operator that assigns to X its picard group, usage \Pic(X)
% \Picardstack_{X/B} denotes the Picard stack of X over B
% \Picardfunctor_{X/B} denotes the Picard functor of X over B
\def\Pic{\mathop{\mathrm{Pic}}\nolimits}
\def\Picardstack{\mathcal{P}\!\mathit{ic}}
\def\Picardfunctor{\mathrm{Pic}}
\def\Deformationcategory{\mathcal{D}\!\mathit{ef}}

\setlength{\emergencystretch}{3em}
\begin{document}
\title[Stacks Project, Tag 07P4]{345 --- Filtered subfield intersections and characteristic-p differentials}
\maketitle
\noindent Copyright (C) 2005--2025 Johan de Jong. Stacks Project authors. Source: \href{https://stacks.math.columbia.edu/tag/07P4}{Tag 07P4}.

\noindent Editorial difficulty note (not author text). Difficulty H2: The author detects a differential relation through a p-basis and uses directed subfield intersections to force its coefficients into pth powers. For finite field extensions it reduces through a tower to separable and purely inseparable primitive cases and proves compatibility of the intersection..

\noindent Editorial provenance note (not author text). History: excerpt from revision \texttt{a0444\allowbreak 6e57e\allowbreak c1fbc\allowbreak 252a8\allowbreak 71afc\allowbreak ec775\allowbreak 2fb28\allowbreak 07b14}, more-algebra.tex, lines 11457--11518. Reference presentation was adapted; original-author context and core support blocks were retained or added. The mathematical wording is unchanged. Licensed under GNU FDL 1.2 or later; no invariant sections or cover texts. See ../licenses/COPYING. Context blocks reproduce author definitions and supporting results; they are not additional counted problems.

\begin{definition}
\label{definition-p-basis}
Let $p$ be a prime number. Let $k \to K$ be an extension of fields
of characteristic $p$. Denote $kK^p$ the compositum of $k$ and $K^p$
in $K$.
\begin{enumerate}
\item A subset $\{x_i\} \subset K$ is called {\it p-independent
over $k$} if the elements $x^E = \prod x_i^{e_i}$ where
$0 \leq e_i < p$ are linearly independent over $kK^p$.
\item A subset $\{x_i\}$ of $K$ is called a
{\it p-basis of $K$ over $k$} if the elements
$x^E$ form a basis of $K$ over $kK^p$.
\end{enumerate}
\end{definition}

\begin{lemma}
\label{lemma-intersection-subfields}
Let $K$ be a field of characteristic $p$. Let $\{K_\alpha\}_{\alpha \in A}$
be a collection of subfields of $K$ with the following properties
\begin{enumerate}
\item $K^p \subset K_\alpha$ for all $\alpha \in A$,
\item $K^p = \bigcap_{\alpha \in A} K_\alpha$,
\item for $\alpha, \alpha' \in A$ there exists an $\alpha'' \in A$
such that $K_{\alpha''} \subset K_\alpha \cap K_{\alpha'}$.
\end{enumerate}
Then
\begin{enumerate}
\item the intersection of the kernels of the maps
$\Omega_{K/\mathbf{F}_p} \to \Omega_{K/K_\alpha}$ is zero,
\item for any finite extension $L/K$ we have
$L^p = \bigcap_{\alpha \in A} L^pK_\alpha$.
\end{enumerate}
\end{lemma}

\begin{proof}
Proof of (1).
Choose a $p$-basis $\{x_i\}$ for $K$ over $\mathbf{F}_p$.
Suppose that $\eta = \sum_{i \in I'} y_i \text{d}x_i$ maps to zero in
$\Omega_{K/K_\alpha}$ for every $\alpha \in A$. Here the index set
$I'$ is finite. By Lemma \href{https://stacks.math.columbia.edu/tag/07P2}{lemma p basis (Tag 07P2)}
this means that for every $\alpha$ there exists a relation
$$
\sum\nolimits_E a_{E, \alpha} x^E = 0,\quad a_{E, \alpha} \in K_\alpha
$$
where $E$ runs over multi-indices $E = (e_i)_{i \in I'}$ with
$0 \leq e_i < p$. On the other hand, Lemma \href{https://stacks.math.columbia.edu/tag/07P2}{lemma p basis (Tag 07P2)}
guarantees there is no such relation $\sum a_E x^E = 0$ with
$a_E \in K^p$. This is a contradiction by
Lemma \href{https://stacks.math.columbia.edu/tag/07P3}{lemma intersection subfields subspace (Tag 07P3)}.

\medskip\noindent
Proof of (2). Suppose that we have a tower $L/M/K$
of finite extensions of fields. Set $M_\alpha = M^p K_\alpha$
and $L_\alpha = L^p K_\alpha = L^p M_\alpha$. Then we can first prove that
$M^p = \bigcap_{\alpha \in A} M_\alpha$, and after that prove
that $L^p = \bigcap_{\alpha \in A} L_\alpha$. Hence it suffices to
prove (2) for primitive field extensions having no nontrivial subfields.
First, assume that $L = K(\theta)$ is separable over $K$. Then
$L$ is generated by $\theta^p$ over $K$, hence we may assume that
$\theta \in L^p$. In this case we see that
$$
L^p = K^p \oplus K^p\theta \oplus \ldots K^p\theta^{d - 1}
\quad\text{and}\quad
L^pK_\alpha =
K_\alpha \oplus K_\alpha \theta \oplus \ldots K_\alpha\theta^{d - 1}
$$
where $d = [L : K]$. Thus the conclusion is clear in this case. The other
case is where $L = K(\theta)$ with $\theta^p = t \in K$, $t \not \in K^p$.
In this case we have
$$
L^p = K^p \oplus K^pt \oplus \ldots K^pt^{p - 1}
\quad\text{and}\quad
L^pK_\alpha =
K_\alpha \oplus K_\alpha t \oplus \ldots K_\alpha t^{p - 1}
$$
Again the result is clear.
\end{proof}
\end{document}
